% --- PACKAGES ---
% Mapeia glifos para Unicode: sem isso, ATS extraem "financial" como "nancial"
% (ligaduras fi/fl) e podem perder acentos.
\input{glyphtounicode}
\pdfgentounicode=1
\usepackage[utf8]{inputenc}
\usepackage[T1]{fontenc}
% Bitstream Charter: desenhada para continuar legível em baixa resolução,
% com traço mais firme que a Latin Modern. O recrutador lê o PDF na tela.
\usepackage{charter}
\usepackage{microtype}
\usepackage{geometry}
\usepackage{parskip}
\usepackage{hyperref}
\usepackage{titlesec}
\usepackage{enumitem}
\usepackage{setspace}

% --- DOCUMENT CONFIG ---
% Margens laterais maiores: linhas de ~100 caracteres em vez de ~120.
\geometry{top=0.9cm, bottom=0.8cm, left=1.4cm, right=1.4cm}
\pagestyle{empty}
\raggedbottom

% --- SPACING ---
% Respiro entre bullets e entre seções: menos "parede de texto".
\setlist[itemize]{itemsep=1pt, topsep=2pt, parsep=0pt, partopsep=0pt, leftmargin=1.1em}
\setlength{\parskip}{1pt}
\setlength{\parindent}{0pt}

\hypersetup{
	pdftitle={Leandro Victtorio Costa Campelo},
	pdfauthor={Leandro Victtorio Costa Campelo},
	colorlinks=true,
	linkcolor=black,
	urlcolor=black
}

\setcounter{secnumdepth}{0}

\titleformat{\section}
{\large\bfseries}{\thesection}{0em}{}
[\vspace{0.3ex}\titlerule]

\titlespacing*{\section}{0pt}{5pt}{3pt}

% Datas à direita; se não couberem, vão inteiras para a linha de baixo,
% ainda alinhadas à direita (truque do \signed, TeXbook p. 106).
\newcommand{\cvdates}[1]{%
  {\unskip\nobreak\hfil\penalty50\hskip1.5em\hbox{}\nobreak\hfil
   \mbox{\small\itshape #1}\parfillskip=0pt\finalhyphendemerits=0\par}}

% Experiência e projetos, em uma linha:
%   \cventry{empresa ou projeto}{cargo}{local (pode ficar vazio)}{datas}
% A empresa vem primeiro e em negrito: é o que o recrutador procura ao bater o olho.
% Tudo na mesma linha de texto, separado por "|": com o local alinhado à direita,
% parsers de ATS liam a cidade como outra empresa.
\newcommand{\cventry}[4]{%
  \par\addvspace{5pt}\noindent
  {\bfseries #1} \textbar{} #2%
  \if\relax\detokenize{#3}\relax\else{} \textbar{} #3\fi
  \cvdates{#4}%
  \nopagebreak[4]}

% Formação: curso em negrito com as datas; instituição na linha de baixo
% (o nome completo da instituição não cabe na mesma linha das datas).
%   \eduentry{curso}{instituição}{datas}
\newcommand{\eduentry}[3]{%
  \par\addvspace{5pt}\noindent
  {\bfseries #1}\cvdates{#3}%
  \nopagebreak[4]\noindent #2\par\nopagebreak[4]}
